\usepackage{physics}

% colored text
\newcommand{\blue}[1]{\textcolor{blue}{#1}}
\newcommand{\red}[1]{\textcolor{red}{#1}}
\newcommand{\green}[1]{\textcolor{green}{#1}}

\usepackage{graphicx} % Для использования команды \resizebox

\makeatletter
\renewcommand\@biblabel[1]{#1.} % Настройка номеров в списке литературы без скобок
\makeatother

\usepackage[absolute,overlay]{textpos}

\usepackage{colortbl} % Для работы с цветами в таблицах
\usepackage{xcolor}   % Для работы с цветами и градиентами

% Определяем диапазон значений
\newcommand{\minval}{1}
\newcommand{\maxval}{22}

\usepackage{float} % Для использования модификатора [H]

\usepackage{amsmath}
\usepackage{booktabs}

\usepackage{afterpage}

% Оформление субсекций в стиле параграфов
\setsubsecheadstyle{\normalfont\normalsize\bfseries}
\setbeforesubsecskip{\baselineskip}
\setaftersubsecskip{-1em}

% Постановка точки в конце заголовков субсекций
\makeatletter
\let\oldsubsection\subsection

\renewcommand{\subsection}{%
    \@ifstar{\subsectionstar}{\subsectionnostar}%
}

\newcommand{\subsectionstar}[1]{%
    \oldsubsection*{#1.}%
}

\newcommand{\subsectionnostar}[1]{%
    \oldsubsection{#1.}%
}
\makeatother

% Круглая точка в маркированных списках
\renewcommand{\labelitemi}{\textbullet}

% Белая линия для \multicolumn
\newcommand{\whitehline}{%
    \arrayrulecolor{gray}\hline\arrayrulecolor{black}%
}

% Библиотека для построения графиков с осью дат
\usepgfplotslibrary{dateplot}

% Корректное отображение ссылок
\usepackage{xurl}